\documentclass[11pt, a4paper]{article}

\usepackage[margin=2.2cm]{geometry}
\usepackage{amsmath, amssymb}
\usepackage{bm}
\usepackage{booktabs}
\usepackage{xcolor}
\usepackage{hyperref}
\usepackage{tcolorbox}
\usepackage{parskip}
\tcbuselibrary{skins}

\hypersetup{colorlinks=true, linkcolor=blue!60!black}

\newcommand{\R}{\mathbb{R}}
\newcommand{\Sig}{\boldsymbol{\Sigma}}
\newcommand{\sighat}{\hat{\boldsymbol{\Sigma}}}
\newcommand{\KL}{\mathrm{KL}}
\newcommand{\T}{{\top}}

\newtcolorbox{eq}[1][]{
    enhanced, sharp corners,
    colback=gray!8, colframe=black!30,
    boxrule=0.6pt, left=6pt, right=6pt, top=4pt, bottom=4pt,
    fontupper=\small, #1
}

\title{\textbf{Soft Anchor-based Knowledge Distillation}\\[0.2em]
\normalsize Methodology Summary}
\author{}
\date{\today}

\begin{document}
\maketitle
\vspace{-1em}

% -------------------------------------------------------
\section*{1\quad Anchor-based Knowledge Distillation (AKD)}
% -------------------------------------------------------

\textbf{Anchors.} One representative image per class is selected from the
training set as the sample with highest cosine centrality among its
class peers (i.e.\ highest average cosine similarity to all other samples
of the same class). This gives an anchor set $\mathcal{A}$ of $C=100$ images.
A learnable per-anchor spatial attention map $\mathbf{W}_c$ (energy-normalised
at each step) is applied to the anchor before feeding it to the student
network (\emph{AnchorNet}).

\textbf{Similarity matrices.} Let $\mathbf{T},\mathbf{S}\in\R^{B\times D}$
be $\ell_2$-normalised teacher/student batch features and
$\mathbf{T}_A,\mathbf{S}_A\in\R^{A\times D}$ the corresponding anchor features.
Cosine similarity is scaled to $[0,1]$ via $s(\mathbf{u},\mathbf{v})=(\mathbf{u}^\T\mathbf{v}+1)/2$.
Three teacher/student matrix pairs are formed:
\[
    \mathbf{P}^{BB} = \frac{\mathbf{T}\mathbf{T}^\T+\mathbf{1}}{2}\in\R^{B\times B},
    \quad
    \mathbf{P}^{BA} = \frac{\mathbf{T}\mathbf{T}_A^\T+\mathbf{1}}{2}\in\R^{B\times A},
    \quad
    \mathbf{P}^{AB} = \bigl(\mathbf{P}^{BA}\bigr)^\T\in\R^{A\times B}.
\]
Each is row-normalised to a probability distribution
$\tilde{\mathbf{P}}_{i\cdot} = \mathbf{P}_{i\cdot}/\sum_j P_{ij}$.

\textbf{Loss.}
\begin{eq}
\[
    \mathcal{L}_k = \sum_i \KL\!\left(\tilde{P}^k_{T,i\cdot}\;\|\;\tilde{P}^k_{S,i\cdot}\right),
    \quad k\in\{BB,\,BA,\,AB\}
    \qquad
    \mathcal{L}_{\text{AKD}} = \ell_1\,\mathcal{L}_{BB}
    + (1-\ell_2)\,\mathcal{L}_{BA}
    + \ell_2\,\mathcal{L}_{AB}
\]
\end{eq}
Total objective: $\mathcal{L} = \mathcal{L}_{\text{CE}} + \beta\,\mathcal{L}_{\text{AKD}}$.

% -------------------------------------------------------
\section*{2\quad Class Similarity Matrix $\Sig$}
% -------------------------------------------------------

A global $C\times C$ semantic prior is computed \emph{once} from the frozen
teacher over the full training set:
\[
    \Sig_{ij} = \frac{1}{|\mathcal{F}_i||\mathcal{F}_j|}
    \sum_{\mathbf{u}\in\mathcal{F}_i}\sum_{\mathbf{v}\in\mathcal{F}_j}
    \frac{\mathbf{u}^\T\mathbf{v}+1}{2},
\]
where $\mathcal{F}_c$ is the set of $\ell_2$-normalised teacher features for class $c$.
$\Sig$ is symmetric and reflects the teacher's inter-class similarity structure.

\textbf{Temperature sharpening} ($\tau < 1$ increases contrast):
\[
    \Sig^{(\tau)}_{ij} = \Sig_{ij}^{1/\tau}.
\]

% -------------------------------------------------------
\section*{3\quad SoftAKD: Blending Instance and Class-level Knowledge}
% -------------------------------------------------------

For the current batch/anchor labels $\mathbf{y}, \mathbf{y}_A$, we extract
label-indexed subsets of $\Sig^{(\tau)}$:
$\;\Sig_{BB}=\Sig^{(\tau)}[\mathbf{y},\mathbf{y}]$,\;
$\Sig_{BA}=\Sig^{(\tau)}[\mathbf{y},\mathbf{y}_A]$,\;
$\Sig_{AB}=\Sig^{(\tau)}[\mathbf{y}_A,\mathbf{y}]$.

Then each teacher similarity matrix is softened before the KL computation.

\subsection*{3a\quad Static $\lambda$ (baseline)}

A single scalar $\lambda\in[0,1]$ blends all rows uniformly:
\begin{eq}
\[
    \mathbf{P}^{k,\text{soft}}_{T}
    = \lambda\,\mathbf{P}^k_{T} + (1-\lambda)\,\Sig_k,
    \qquad k\in\{BB,\,BA,\,AB\}.
\]
\end{eq}
At $\lambda=1$ this reduces to standard AKD; at $\lambda=0.9$ (used in experiments)
10\% of the class-level structure is injected into every sample equally.

\subsection*{3b\quad Adaptive $\lambda$ via Sigmoid Gate}

The scalar $\lambda$ cannot distinguish easy samples (student already agrees
with teacher) from hard ones. We replace it with a \emph{per-row} weight
that increases class-level blending only when the student's distribution
disagrees with the teacher.

\textbf{Step 1 — Pearson distance} between row $i$ of the teacher and
student similarity matrices:
\begin{eq}
\[
    d_i = 1 - \frac{(\mathbf{p}_i^T - \bar{p}_i)^\T(\mathbf{p}_i^S - \bar{p}_i^S)}
    {\|\mathbf{p}_i^T-\bar{p}_i^T\|_2\,\|\mathbf{p}_i^S-\bar{p}_i^S\|_2}
    \;\in[0,2].
\]
\end{eq}
Mean-centring removes absolute level effects; $d_i$ measures shape disagreement.

\textbf{Step 2 — Normalise and apply sigmoid gate:}
\begin{eq}
\[
    \tilde{d}_i = \frac{d_i}{\max_j d_j},
    \qquad
    g_i = \sigma\!\left(k\,(\tilde{d}_i - d_0)\right),
    \qquad
    \lambda_i = 1 - g_i\,(1-\lambda_{\min})
    \;\in[\lambda_{\min},\,1].
\]
\end{eq}

\textbf{Step 3 — Per-row blending} (applied independently to each of $BB$, $BA$, $AB$):
\begin{eq}
\[
    \mathbf{P}^{k,\text{soft}}_{T,i\cdot}
    = \lambda_i^k\,\mathbf{P}^k_{T,i\cdot}
    + (1-\lambda_i^k)\,\Sig_{k,i\cdot}.
\]
\end{eq}

\begin{center}
\small
\begin{tabular}{lll}
\toprule
 & Static ($\lambda=0.9$) & Adaptive \\
\midrule
Granularity & Same for all samples & Per sample, per term, per step \\
Range & Fixed scalar & $[\lambda_{\min},\,1.0]$ per row \\
Responds to disagreement & No & Yes \\
Extra hyperparameters & None & $\lambda_{\min}=0.7$,\; $k=10$,\; $d_0=0.3$ \\
\bottomrule
\end{tabular}
\end{center}

% -------------------------------------------------------
\section*{4\quad CASL: Class-Adaptive Softening of Labels \cite{spanos}}
% -------------------------------------------------------

CASL addresses class uncertainty by learning a $K\times K$ confusion matrix
$\mathbf{M}$ that redistributes probability mass from a one-hot label toward
semantically similar classes.

\textbf{Initialization.}
$\mathbf{M}$ is initialised as the identity plus small Gaussian noise,
$\mathbf{M}\leftarrow\mathbf{I}+\boldsymbol{\varepsilon}$
($\varepsilon_{ij}\sim\mathcal{N}(0,10^{-4})$),
so training begins near one-hot labels with diagonal dominance guaranteed from the start.

\textbf{Soft label generation.}
For a one-hot label $\mathbf{y}\in\{0,1\}^K$, the CASL soft target is:
\begin{eq}
\[
    g(\mathbf{y}) = \mathrm{softmax}(\mathbf{M}\mathbf{y}),
    \qquad
    \hat{\mathbf{y}}_{\text{CASL}} = \alpha\,g(\mathbf{y}) + (1-\alpha)\,\mathbf{y},\quad \alpha=0.1.
\]
\end{eq}
Since $\mathbf{y}$ is one-hot for class $c$, $\mathbf{M}\mathbf{y}$ selects
column $c$ of $\mathbf{M}$; the softmax normalises it into a valid distribution
over all classes. $\alpha$ controls softness while preserving the ground-truth signal.

\textbf{Dual loss and training.}
Two losses are jointly minimised at every step:
\begin{eq}
\[
    \mathcal{L}_g = -\frac{1}{K}\sum_k \bigl[\phi(\mathbf{x})\bigr]_k \log\bigl[g(\mathbf{y})\bigr]_k
                   + (1-[\phi]_k)\log(1-[g]_k),
    \qquad
    \mathcal{L}_{\text{CASL}} = \mathcal{L}_g + \mathcal{L}(\mathbf{x},\hat{\mathbf{y}}_{\text{CASL}}).
\]
\end{eq}
$\mathcal{L}_g$ uses binary cross-entropy with the \emph{classifier's own output}
$\phi(\mathbf{x})$ as the supervision target for $\mathbf{M}$, forcing $\mathbf{M}$ to
encode the confusion structure the classifier itself discovers.
$\mathcal{L}(\mathbf{x},\hat{\mathbf{y}}_{\text{CASL}})$ then trains the classifier
with the resulting soft targets.
The two objectives create a \emph{cooperative loop}: the classifier's predictions
supervise $\mathbf{M}$, while $\mathbf{M}$'s soft labels regularise the classifier.
$\mathbf{M}$ is optimised with a separate Adam optimiser (lr $10^{-5}$ on CIFAR-100)
and adds negligible overhead ($<\!2\%$ extra training time; $K^2=10{,}000$ parameters
for $K=100$).

% -------------------------------------------------------
\section*{5\quad GCN: Adapting $\Sig$ Inspired by CASL}
% -------------------------------------------------------

The static $\Sig$ does not adapt to the student's evolving confusion.
Inspired by CASL \cite{spanos}, we introduce a learnable $C\times C$
matrix $\mathbf{M}$ (initialised as $\mathbf{I}+\boldsymbol{\varepsilon}$,
$\varepsilon\sim\mathcal{N}(0,10^{-4})$) that adjusts $\Sig$ to track the
student's current inter-class similarity pattern.

\textbf{Adapted sigma:}
\begin{eq}
\[
    \sighat = \alpha\,(\Sig\,\mathbf{M}) + (1-\alpha)\,\Sig.
\]
\end{eq}

\textbf{Student confusion target $\Sig_S$.}
At each step, the student's pairwise similarities are aggregated into a
class-level matrix:
\[
    (\Sig_S)_{mn} = \frac{1}{|\{(i,j):c_i=m,c_j=n\}|}
    \sum_{i:c_i=m}\sum_{j:c_j=n}
    \frac{\mathbf{s}_i^\T\mathbf{s}_j+1}{2}.
\]
Three modes control which samples form rows/columns:
\texttt{ab} (batch $\times$ anchors, default), \texttt{aa} (anchors $\times$ anchors),
\texttt{bb} (batch $\times$ batch).

\textbf{GCN loss} (optimised with a separate Adam optimiser; $\sighat$ is
detached from the student graph):
\begin{eq}
\[
    \mathcal{L}_G
    = \mathrm{MSE}\!\left(\sighat\odot\mathbf{M}_{\text{obs}},\;
                          \Sig_S\odot\mathbf{M}_{\text{obs}}\right),
\]
\end{eq}
where $\mathbf{M}_{\text{obs}}$ masks out class pairs not observed in the batch.
$\sighat$ then replaces $\Sig$ in all subsequent blending steps.

% -------------------------------------------------------
\section*{6\quad Hyperparameters}
% -------------------------------------------------------

\begin{center}
\small
\begin{tabular}{llll}
\toprule
Symbol & Meaning & Default \\
\midrule
$\beta$ & AKD distillation weight & 10.0 \\
$\ell_1,\,\ell_2$ & KL term weights & 1.0,\;0.1 \\
$\tau$ & Sigma temperature ($<1$ sharpens) & 0.2 \\
$\alpha$ & GCN blending weight & 0.1 \\
$\eta_G$ & GCN learning rate & $10^{-4}$ \\
$\lambda_{\min}$ & Adaptive lambda floor & 0.7 \\
$k$ & Sigmoid steepness & 10.0 \\
$d_0$ & Sigmoid midpoint (normalised disagreement) & 0.3 \\
\bottomrule
\end{tabular}
\end{center}

\vspace{0.5em}
\begin{thebibliography}{1}
\bibitem{spanos}
G.~Spanos et al., \textit{Improving uncertainty estimation in deep neural networks by modeling class and instance uncertainty}, 2019.
\end{thebibliography}

\end{document}
